\section{Generating Set for the LPS Family}
  \label{sec:step6a}

  The spectral admissibility programme follows the chain
  \begin{equation}
    Q_8
    \;\longrightarrow\;
    2I
    \;\longrightarrow\;
    \mathrm{SL}(2,\mathbb{F}_p)
    \;\xrightarrow{\;p \to \infty\;}
    \mathrm{SU}(2),
    \label{eq:chain-6a}
  \end{equation}
  where the final arrow denotes local spectral convergence toward the
  $S^3$ Laplacian.
  The first concrete step of the third stage is an explicit symmetric
  generating set whose structural role is analogous to the order-4
  generators of $Q_8$ and the order-10 generators of $2I$.
  This section constructs such a set via the Lubotzky--Phillips--Sarnak
  (LPS) construction~\cite{LPS1988}.

% ------------------------------------------------------------
  \subsection{The groups and their sizes}
    \label{sec:step6a-group}

    The LPS construction involves two primes $p$ and $q$, both congruent to
    $1 \pmod{4}$, with $p \neq q$.
    The Cayley graph $X^{p,q}$ has vertex set $\mathrm{PSL}(2,\mathbb{F}_q)$
    and degree $p + 1$.
    Its order is
    \begin{equation}
      |\mathrm{PSL}(2,\mathbb{F}_q)| = \tfrac{1}{2}\,q(q^2-1).
      \label{eq:PSL2-order}
    \end{equation}
    As $q \to \infty$ with $p$ fixed, the group grows without bound and the
    Cayley graph approximates the geometry of $S^3$ progressively more finely.

    \medskip
    \begin{center}
      \renewcommand{\arraystretch}{1.2}
      \begin{tabular}{ccc}
        \hline
        $p$ & $q$ & $|\mathrm{PSL}(2,\mathbb{F}_q)|$ \\
        \hline
        5 & 29 & 12180 \\
        5 & 41 & 34440 \\
        13 & 17 & 2448  \\
        13 & 29 & 12180 \\
        \hline
      \end{tabular}
    \end{center}

    \medskip\noindent
    The role of the two primes is asymmetric: $p$ determines the degree
    $d = p+1$ (and hence the generating set structure), while $q$ determines
    the size of the group.
    Step~6 fixes $p$ and varies $q \to \infty$.

% ------------------------------------------------------------
  \subsection{Quaternion algebra and the spinor embedding}
    \label{sec:step6a-quat}

    The Hurwitz quaternion algebra $\mathbb{H}(\mathbb{Z})$ embeds into
    $M_2(\mathbb{F}_q)$ via the standard spinor map.
    Fix a square root of $-1$ in $\mathbb{F}_q$, denoted $\mathbf{i}_q$
    (which exists since $q \equiv 1 \pmod{4}$).
    The map
    \begin{equation}
      \Phi_q \colon \mathbb{H}(\mathbb{Z}) \longrightarrow M_2(\mathbb{F}_q),
      \qquad
      a + b\mathbf{i} + c\mathbf{j} + d\mathbf{k}
      \;\longmapsto\;
      \begin{pmatrix}
        a + b\,\mathbf{i}_q & c + d\,\mathbf{i}_q \\
        -c + d\,\mathbf{i}_q & a - b\,\mathbf{i}_q
      \end{pmatrix}
      \pmod{q}
      \label{eq:spinor-map}
    \end{equation}
    sends a quaternion of integer norm $a^2+b^2+c^2+d^2$ to a matrix of the
    same determinant modulo $q$.
    For a quaternion of norm $p$, one has
    $\det \Phi_q(\cdot) \equiv p \pmod{q}$.
    To normalise the determinant to $1$, one multiplies by $r^{-1} \in
\mathbb{F}_q$ where $r$ satisfies $r^2 \equiv p \pmod{q}$.
    Such $r$ exists provided $p$ is a quadratic residue modulo $q$.

    \begin{remark}
      The normalisation factor $r^{-1}$ is an element of $\mathbb{F}_q$,
      not of $\mathbb{F}_p$.
      There is no division by $\sqrt{p}$ in $\mathbb{F}_p$; the entire
      construction takes place in $\mathbb{F}_q$.
    \end{remark}

% ------------------------------------------------------------
  \subsection{The LPS generating set}
    \label{sec:step6a-lps}

    \begin{definition}[LPS generators~\cite{LPS1988}]
      \label{def:lps}
      Let $p, q$ be distinct primes with $p \equiv q \equiv 1 \pmod{4}$ and
      $\bigl(\frac{p}{q}\bigr) = 1$.
      Fix $\mathbf{i}_q$ with $\mathbf{i}_q^2 \equiv -1 \pmod{q}$, and
      $r \in \mathbb{F}_q$ with $r^2 \equiv p \pmod{q}$.
      Let $\Omega_p$ be the set of integer solutions to
      \begin{equation}
        a^2 + b^2 + c^2 + d^2 = p,
        \qquad a > 0,\; a \text{ odd},\; b, c, d \text{ even}.
        \label{eq:sum-of-4-squares}
      \end{equation}
      By Jacobi's four-square theorem, after quotienting by sign and permutation
      symmetries, there are exactly $p+1$ such solutions.
      For each $(a,b,c,d) \in \Omega_p$, define
      \begin{equation}
        g_{(a,b,c,d)}
        = r^{-1}
        \begin{pmatrix}
          a + b\,\mathbf{i}_q & c + d\,\mathbf{i}_q \\
          -c + d\,\mathbf{i}_q & a - b\,\mathbf{i}_q
        \end{pmatrix}
        \pmod{q}
        \;\in\; \mathrm{SL}(2,\mathbb{F}_q).
        \label{eq:lps-generator}
      \end{equation}
      The LPS generating set is
      \begin{equation}
        S_p = \bigl\{ \pm\, g_{(a,b,c,d)} \mid (a,b,c,d) \in \Omega_p \bigr\}
        \big/\{\pm I\}
        \;\subset\; \mathrm{PSL}(2,\mathbb{F}_q),
        \qquad |S_p| = p+1.
        \label{eq:S-SL2}
      \end{equation}
    \end{definition}

    The set $S_p$ is symmetric ($g \in S_p \Rightarrow g^{-1} \in S_p$)
    because $(-a,-b,-c,-d) \in \Omega_p$ whenever $(a,b,c,d) \in \Omega_p$,
    and $g_{(-a,-b,-c,-d)} = g_{(a,b,c,d)}^{-1}$.
    The Cayley graph $X^{p,q} = \mathrm{Cay}(\mathrm{PSL}(2,\mathbb{F}_q), S_p)$
    is therefore undirected and the associated Laplacian self-adjoint.

% ------------------------------------------------------------
  \subsection{Structural analogy with $Q_8$ and $2I$}
    \label{sec:step6a-analogy}

    The structural role of $S_p$ is the exact analogue of the generating sets
    used on $Q_8$ and $2I$.
    Table~\ref{tab:gensets} summarises the three parallel constructions.

    \begin{table}[h]
      \centering
      \renewcommand{\arraystretch}{1.3}
      \caption{Parallel structure of generating sets across the programme chain.}
      \label{tab:gensets}
      \begin{tabular}{llll}
        \hline
        Group & $|G|$ & Generator type & $|S|$ \\
        \hline
        $Q_8$ & $8$ &
        order-4 elements ($\pm i,\pm j,\pm k$);
        spinorial lifts of $\pi$-rotations in $SO(3)$ & $6$ \\
        $2I \cong \mathrm{SL}(2,\mathbb{F}_5)$ & $120$ &
        order-10 elements ($10a \cup 10b$);
        spinorial lifts of order-5 rotations in $A_5$ & $6$ \\
        $\mathrm{PSL}(2,\mathbb{F}_q)$ or $\mathrm{PGL}(2,\mathbb{F}_q)$ &
        $\tfrac{1}{2}q(q^2{-}1)$ or $q(q^2{-}1)$ &
        LPS generators from norm-$p$ quaternions
        via spinor map into $\mathbb{F}_q$ & $p{+}1$ \\
        \hline
      \end{tabular}
    \end{table}

    In each case the generating set arises from lifting rotational generators
    of the quotient ($SO(3)$, $A_5$, or $\mathrm{PGL}(2,\mathbb{F}_q)$) to
    the spinorial cover via the quaternion map.
    The precise target group is $\mathrm{PSL}(2,\mathbb{F}_q)$ or
    $\mathrm{PGL}(2,\mathbb{F}_q)$ depending on the arithmetic conditions on
    the pair $(p,q)$~\cite{LPS1988}; both cases yield Ramanujan graphs with
    the same spectral bound~\eqref{eq:ramanujan-bound}.
    The degree $d = p+1$ is fixed by $p$ independently of $q$; the group size
    is controlled by $q$.
    Step~6 exploits this decoupling: fix $p$ and vary $q \to \infty$.

% ------------------------------------------------------------
  \subsection{The Ramanujan property}
    \label{sec:step6a-ramanujan}

    The graph $X^{p,q}$ is a $(p+1)$-regular Ramanujan graph~\cite{LPS1988}.
    This means that all non-trivial adjacency eigenvalues $\mu$
    (i.e.\ all except the trivial eigenvalue $\mu = p+1$) satisfy
    \begin{equation}
      |\mu| \leq 2\sqrt{p},
      \label{eq:ramanujan-bound}
    \end{equation}
    which saturates the Alon--Boppana lower bound on the spectral gap.
    Equivalently, all non-trivial Laplacian eigenvalues satisfy
    \begin{equation}
      \lambda = (p+1) - \mu \geq \bigl(\sqrt{p} - 1\bigr)^2 > 0.
      \label{eq:ramanujan-gap}
    \end{equation}
    The admissibility envelope $A^{\max}_\rho = c_\chi/\sqrt{\lambda_\rho}$
    therefore remains well-defined and bounded away from infinity for every
    non-trivial representation $\rho$, uniformly for all $q$.

    The proof uses the Ramanujan--Petersson conjecture for holomorphic cusp
    forms of weight~2, established by Eichler and Shimura.
    The sum-of-four-squares condition~\eqref{eq:sum-of-4-squares} is what
    connects the adjacency operator of $X^{p,q}$ to Hecke operators on
    modular forms, whose eigenvalues are controlled by the Ramanujan bound.

% ------------------------------------------------------------
  \subsection{Explicit example: $p = 5$, $q = 29$}
    \label{sec:step6a-p5q29}

    The case $p = 5$ is structurally connected to Step~5.
    The $p+1 = 6$ solutions to~\eqref{eq:sum-of-4-squares} are
    \begin{equation}
      \Omega_5 = \bigl\{
      (1, \pm 2, 0, 0),\;
      (1, 0, \pm 2, 0),\;
      (1, 0, 0, \pm 2)
      \bigr\},
      \label{eq:omega5}
    \end{equation}
    which are the same norm-5 integer quaternions that, when lifted to
    $2I \cong \mathrm{SL}(2,\mathbb{F}_5)$ via the spinor map, give the
    order-10 generating set $10a \cup 10b$ studied in Step~5.
    The structural analogy is therefore exact at the level of the integer
    quaternion data $\Omega_p$; the target groups differ ($\mathrm{SL}(2,
\mathbb{F}_5)$ in Step~5 versus $\mathrm{PSL}(2,\mathbb{F}_q)$ or
    $\mathrm{PGL}(2,\mathbb{F}_q)$ with $q > 5$ in Step~6), but the
    generating mechanism is the same.

    For $q = 29$: one has $\mathbf{i}_{29} = 12$ (since $12^2 \equiv -1
\pmod{29}$) and $r = 11$ (since $11^2 \equiv 5 \pmod{29}$), giving
    $r^{-1} = 8$.
    Three of the six generators in $\mathrm{SL}(2,\mathbb{F}_{29})$ are
    \begin{equation}
      g_{(1,2,0,0)} =
      \begin{pmatrix} 26 & 0 \\ 0 & 19 \end{pmatrix}, \quad
      g_{(1,0,2,0)} =
      \begin{pmatrix} 8 & 16 \\ 13 & 8 \end{pmatrix}, \quad
      g_{(1,0,0,2)} =
      \begin{pmatrix} 8 & 18 \\ 18 & 8 \end{pmatrix},
      \label{eq:generators-p5q29}
    \end{equation}
    each with $\det = 1$ verified.
    The remaining three generators are the inverses of these.
    The resulting graph $X^{5,29}$ is a 6-regular Ramanujan graph on
    $12180$ vertices.

% ------------------------------------------------------------
  \subsection{The normalised character correlation: key observable}
    \label{sec:step6a-coherence}

    With $S_p$ defined, the key observable for each low-dimensional
    representation sector $\rho$ of $\mathrm{PSL}(2,\mathbb{F}_q)$ is the
    \emph{normalised character correlation}
    \begin{equation}
      \kappa_\rho \;\equiv\; \frac{\mu_\rho}{p+1}
      = \frac{1}{(p+1)\dim\rho} \sum_{g \in S_p} \chi_\rho(g).
      \label{eq:coherence}
    \end{equation}
    As a guiding heuristic, the low-dimensional representations of
    $\mathrm{PSL}(2,\mathbb{F}_q)$ play the role of discrete analogues of
    the low-spin representations of $\mathrm{SU}(2)$ in the large-$q$ limit.
    We index them by $j$ accordingly, writing $\kappa_j$ for $\kappa_{\rho_j}$,
    with the understanding that the identification with $\mathrm{SU}(2)$ spins
    is an asymptotic guide, not a formal equivalence at finite $q$.
    The quantity $\kappa_j$ measures the coherence between $S_p$ and the
    representation $\rho_j$ per generator per dimension.

    The results of Steps~4 and~5 already show that $\kappa_j$ decreases
    with $j$ in the two smallest groups:
    \begin{itemize}
      \item On $Q_8$: $\kappa_{1/2} = 0$ and $\kappa_1 = -1/3 < 0$
      (Pauli tracelessness forces $\mu_{\rho_4} = 0$ exactly).
      \item On $2I$: $\kappa_{1/2} > \kappa_1 > \kappa_{3/2} \geq \kappa_2$
      (the golden ratio identity controls the two lowest levels, and the
      hierarchy extends upward).
    \end{itemize}
    The structural reason is the behaviour of the $\mathrm{SU}(2)$ characters
    \begin{equation}
      \chi_j(\theta) = \frac{\sin\bigl((2j+1)\theta/2\bigr)}{\sin(\theta/2)},
      \label{eq:SU2-character}
    \end{equation}
    which oscillate increasingly rapidly as $j$ grows.
    For the LPS generators, which are concentrated at a definite conjugacy
    class in $\mathrm{PSL}(2,\mathbb{F}_q)$, the sum $\sum_{g \in S_p}
\chi_j(g)$ experiences growing cancellations as $j$ increases, because
    the numerator in~\eqref{eq:SU2-character} oscillates while the
    denominator does not.
    The normalised correlation $\kappa_j$ therefore tends to decrease with
    $j$, pushing $\lambda_j$ upward and the admissibility window
    $A^{\max}_j$ downward for large spin.

    The central prediction of Step~6 is that this hierarchy extends to all
    members of the LPS family: for every $p$, the spin-$\tfrac{1}{2}$ sector
    remains the most admissible.
    This is established numerically in Steps~6b--6d and summarised in
    Section~\ref{sec:step6bcd}.

% ------------------------------------------------------------
  \subsection{Summary}
    \label{sec:step6a-summary}

    The LPS generating set $S_p$ of Definition~\ref{def:lps} has the
    following properties:
    \begin{enumerate}
      \item \emph{Explicit.} Defined by the four-square decomposition of $p$
      over $\mathbb{Z}$, embedded in $\mathrm{PSL}(2,\mathbb{F}_q)$ via the
      spinor map with normalisation by $r^{-1} \in \mathbb{F}_q$.
      No free parameters.
      \item \emph{Symmetric.} The Cayley graph is undirected; the Laplacian
      is self-adjoint.
      \item \emph{Ramanujan.} The non-trivial spectral gap
      satisfies~\eqref{eq:ramanujan-gap} exactly, ensuring the
      admissibility envelope is well-controlled for all $q$.
      \item \emph{Structurally analogous.} $S_p$ is the natural continuation
      of the order-4 generators of $Q_8$ and the order-10 generators of
      $2I$ along the chain~\eqref{eq:chain-6a}.
      \item \emph{Structurally analogous to Step~5.}
      The $p = 5$ integer quaternion data $\Omega_5$ are the same norm-5
      quaternions used in $2I \cong \mathrm{SL}(2,\mathbb{F}_5)$.
      Step~5 and Step~6 share the same generating mechanism; they differ
      only in the target group ($\mathrm{SL}(2,\mathbb{F}_5)$ versus the
      LPS family with $q > p$).
    \end{enumerate}
    This construction provides the natural large-group continuation of the
    $Q_8$ and $2I$ cases.
    The resulting Ramanujan Cayley graphs on $\mathrm{PSL}(2,\mathbb{F}_q)$
    or $\mathrm{PGL}(2,\mathbb{F}_q)$ constitute the canonical testing ground
    for the spectral admissibility programme in the limit approaching
    $\mathrm{SU}(2)$, with the normalised character correlation
    $\kappa_\rho$~\eqref{eq:coherence} as the key observable of Step~6.
